\documentclass[12pt,a4paper]{article}

\usepackage[utf8]{inputenc}
\usepackage[T1]{fontenc}
\usepackage[brazil]{babel}
\usepackage[top=3cm,left=3cm,bottom=2cm,right=2cm]{geometry}
\usepackage{mathptmx}
\usepackage{inconsolata}
\usepackage{setspace}
\usepackage{indentfirst}
\usepackage{titlesec}
\usepackage{booktabs}
\usepackage{array}
\usepackage{listings}
\usepackage{xcolor}
\usepackage{amsmath}
\usepackage{float}

\onehalfspacing
\setlength{\parindent}{1.25cm}

\titleformat*{\section}{\normalfont\bfseries\normalsize}
\titleformat*{\subsection}{\normalfont\bfseries\normalsize}
\titleformat*{\subsubsection}{\normalfont\bfseries\normalsize}

\lstset{
    basicstyle=\ttfamily\footnotesize,
    breaklines=true,
    columns=fullflexible,
    keepspaces=true,
    frame=single,
    rulecolor=\color{black!25},
    backgroundcolor=\color{black!4},
    xleftmargin=0.2cm,
    xrightmargin=0.2cm,
    framexleftmargin=4pt,
    framexrightmargin=4pt,
    framextopmargin=4pt,
    framexbottommargin=4pt,
    literate=
        {á}{{\'a}}1 {é}{{\'e}}1 {í}{{\'i}}1 {ó}{{\'o}}1 {ú}{{\'u}}1
        {Á}{{\'A}}1 {É}{{\'E}}1 {Í}{{\'I}}1 {Ó}{{\'O}}1 {Ú}{{\'U}}1
        {ã}{{\~a}}1 {õ}{{\~o}}1 {Ã}{{\~A}}1 {Õ}{{\~O}}1
        {â}{{\^a}}1 {ê}{{\^e}}1 {ô}{{\^o}}1
        {à}{{\`a}}1 {ç}{{\c{c}}}1 {Ç}{{\c{C}}}1
}

\begin{document}

\section{Questão 02: K-vizinhos mais próximos (KNN)}

\subsection{Base de dados}

A base de risco de crédito utilizada na primeira lista contém 30 instâncias. O KNN compara seis atributos preditivos: história de crédito, dívida, garantia, faixa de renda, emprego e faixa etária. Exemplo é o identificador de cada linha e Risco é a classe a prever. Há 11 casos de risco Alto, 9 de risco Moderado e 10 de risco Baixo.

\begin{table}[H]
    \centering
    \caption{Instâncias utilizadas na classificação}
    \setlength{\tabcolsep}{1.5pt}
    \renewcommand{\arraystretch}{0.92}
    \begin{tabular}{>{\raggedright\arraybackslash}p{1.60cm} >{\raggedright\arraybackslash}p{2.35cm} >{\raggedright\arraybackslash}p{1.25cm} >{\raggedright\arraybackslash}p{1.70cm} >{\raggedright\arraybackslash}p{2.65cm} >{\raggedright\arraybackslash}p{2.55cm} >{\raggedright\arraybackslash}p{1.20cm} >{\raggedright\arraybackslash}p{1.75cm}}
        \toprule
        \textbf{Exemplo} & \textbf{História de crédito} & \textbf{Dívida} & \textbf{Garantia} & \textbf{Renda} & \textbf{Emprego} & \textbf{Idade} & \textbf{Risco} \\
        \midrule
        E1 & Ruim & Alta & Nenhuma & \$0 a \$15k & Fixo & Adulto & Alto \\
        E2 & Desconhecida & Alta & Nenhuma & \$15k a \$35k & Autônomo & Jovem & Alto \\
        E3 & Desconhecida & Baixa & Nenhuma & \$15k a \$35k & Fixo & Jovem & Moderado \\
        E4 & Desconhecida & Baixa & Nenhuma & \$0 a \$15k & Desempregado & Adulto & Alto \\
        E5 & Desconhecida & Baixa & Nenhuma & Acima de \$35k & Fixo & Adulto & Baixo \\
        E6 & Desconhecida & Baixa & Adequada & Acima de \$35k & Fixo & Adulto & Baixo \\
        E7 & Ruim & Baixa & Nenhuma & \$0 a \$15k & Autônomo & Jovem & Alto \\
        E8 & Ruim & Baixa & Adequada & Acima de \$35k & Fixo & Adulto & Moderado \\
        E9 & Boa & Baixa & Nenhuma & Acima de \$35k & Fixo & Jovem & Baixo \\
        E10 & Boa & Alta & Adequada & Acima de \$35k & Fixo & Adulto & Baixo \\
        E11 & Boa & Alta & Nenhuma & \$0 a \$15k & Autônomo & Jovem & Alto \\
        E12 & Boa & Alta & Nenhuma & \$15k a \$35k & Fixo & Jovem & Moderado \\
        E13 & Boa & Alta & Nenhuma & Acima de \$35k & Fixo & Adulto & Baixo \\
        E14 & Ruim & Alta & Nenhuma & \$15k a \$35k & Desempregado & Adulto & Alto \\
        E15 & Boa & Baixa & Adequada & \$15k a \$35k & Fixo & Adulto & Baixo \\
        E16 & Ruim & Alta & Nenhuma & \$0 a \$15k & Desempregado & Jovem & Alto \\
        E17 & Desconhecida & Baixa & Adequada & \$15k a \$35k & Autônomo & Adulto & Moderado \\
        E18 & Boa & Baixa & Nenhuma & Acima de \$35k & Fixo & Adulto & Baixo \\
        E19 & Ruim & Alta & Adequada & \$15k a \$35k & Fixo & Adulto & Moderado \\
        E20 & Desconhecida & Alta & Nenhuma & \$0 a \$15k & Desempregado & Adulto & Alto \\
        E21 & Boa & Alta & Adequada & Acima de \$35k & Autônomo & Adulto & Baixo \\
        E22 & Ruim & Baixa & Nenhuma & \$15k a \$35k & Fixo & Jovem & Moderado \\
        E23 & Desconhecida & Baixa & Nenhuma & Acima de \$35k & Fixo & Adulto & Baixo \\
        E24 & Ruim & Alta & Nenhuma & \$0 a \$15k & Autônomo & Jovem & Alto \\
        E25 & Boa & Baixa & Adequada & \$0 a \$15k & Fixo & Adulto & Moderado \\
        E26 & Desconhecida & Alta & Adequada & Acima de \$35k & Fixo & Jovem & Baixo \\
        E27 & Boa & Baixa & Nenhuma & \$15k a \$35k & Autônomo & Jovem & Moderado \\
        E28 & Ruim & Alta & Nenhuma & Acima de \$35k & Desempregado & Adulto & Alto \\
        E29 & Desconhecida & Baixa & Adequada & \$0 a \$15k & Fixo & Adulto & Moderado \\
        E30 & Boa & Alta & Nenhuma & \$0 a \$15k & Fixo & Jovem & Alto \\
        \bottomrule
    \end{tabular}
\end{table}

\subsection{Funcionamento do algoritmo}

O KNN armazena os exemplos rotulados e, a cada consulta, calcula a distância entre o perfil recebido e os perfis da base de referência. A distância de Hamming soma 1 para cada atributo diferente e 0 para cada coincidência. Os seis atributos têm o mesmo peso e a distância varia de 0 a 6.

\[
d(x, y) = \text{número de atributos em que } x \text{ e } y \text{ diferem}
\]

Os registros são ordenados pela distância. Quando dois registros ficam à mesma distância, vem primeiro o de menor número no identificador. Os $k$ primeiros vizinhos votam, cada um com um voto para sua classe. Em caso de empate na votação, vence a classe cujo primeiro representante aparece mais cedo na lista ordenada.

\subsection{Classificação manual de E1 com k = 5}

E1 tem história Ruim, dívida Alta, garantia Nenhuma, renda de \$0 a \$15k, emprego Fixo e idade Adulto. Retiramos E1 da base de referência e comparamos seu perfil aos outros 29. E13 difere na história de crédito e na renda, portanto $d(\text{E1}, \text{E13}) = 1 + 0 + 0 + 1 + 0 + 0 = 2$.

\begin{table}[H]
    \centering
    \caption{Cinco vizinhos selecionados pelo critério de desempate}
    \begin{tabular}{>{\raggedright\arraybackslash}p{1.4cm} >{\raggedright\arraybackslash}p{1.6cm} >{\raggedright\arraybackslash}p{6.0cm} >{\raggedright\arraybackslash}p{1.8cm} >{\raggedright\arraybackslash}p{2.0cm}}
        \toprule
        \textbf{Ordem} & \textbf{Vizinho} & \textbf{Atributos diferentes de E1} & \textbf{Distância} & \textbf{Risco} \\
        \midrule
        1 & E13 & História de crédito e renda & 2 & Baixo \\
        \addlinespace
        2 & E14 & Renda e emprego & 2 & Alto \\
        \addlinespace
        3 & E16 & Emprego e idade & 2 & Alto \\
        \addlinespace
        4 & E19 & Garantia e renda & 2 & Moderado \\
        \addlinespace
        5 & E20 & História de crédito e emprego & 2 & Alto \\
        \bottomrule
    \end{tabular}
\end{table}

Oito registros ficaram à distância 2. Os cinco selecionados deram 3 votos para Alto, 1 para Moderado e 1 para Baixo. A previsão foi \textbf{Alto}, igual à classe registrada para E1.

\subsection{Validação nas 30 instâncias}

A validação \textit{leave-\allowbreak one-\allowbreak out} retira uma instância por rodada e usa as outras 29 para classifi\-cá-la. Após 30 rodadas para cada valor de $k$, a acurácia é a quantidade de acertos dividida por 30.

\begin{table}[H]
    \centering
    \begin{tabular}{>{\raggedright\arraybackslash}p{1.5cm} >{\raggedright\arraybackslash}p{2.8cm} >{\raggedright\arraybackslash}p{2.8cm}}
        \toprule
        \textbf{k} & \textbf{Acertos} & \textbf{Acurácia} \\
        \midrule
        1 & 20 de 30 & 66,7\% \\
        \addlinespace
        3 & 20 de 30 & 66,7\% \\
        \addlinespace
        5 & 21 de 30 & 70,0\% \\
        \addlinespace
        7 & 21 de 30 & 70,0\% \\
        \bottomrule
    \end{tabular}
\end{table}

\subsubsection{Matriz de confusão e análise}

Na matriz de confusão para $k = 5$, as linhas representam as classes registradas no CSV e as colunas representam as classes previstas. A diagonal contém os acertos.

\begin{table}[H]
    \centering
    \begin{tabular}{>{\raggedright\arraybackslash}p{3.2cm} >{\raggedright\arraybackslash}p{2.2cm} >{\raggedright\arraybackslash}p{2.2cm} >{\raggedright\arraybackslash}p{2.2cm}}
        \toprule
        \textbf{Real / Previsto} & \textbf{Alto} & \textbf{Moderado} & \textbf{Baixo} \\
        \midrule
        Alto & 10 & 1 & 0 \\
        \addlinespace
        Moderado & 3 & 2 & 4 \\
        \addlinespace
        Baixo & 0 & 1 & 9 \\
        \bottomrule
    \end{tabular}
\end{table}

O KNN acertou 10 dos 11 casos de risco Alto e 9 dos 10 de risco Baixo. A classe Moderado apresentou o menor número de acertos: 2 de 9. Dos sete erros nessa classe, três casos foram previstos como Alto e quatro como Baixo.

Os valores de $k$ foram comparados na mesma base, por isso os resultados são exploratórios. E5 e E23 têm os mesmos seis atributos. Quando um desses registros é retirado, o outro continua disponível como vizinho a distância zero. A acurácia também depende da regra usada para desempatar distâncias iguais.

\subsection{Execução do programa em Python}

O arquivo \textbf{knn.py} usa a biblioteca padrão do Python e lê o arquivo \textbf{arquivo.csv} da mesma pasta. O comando abaixo demonstra a classificação de E1 com cinco vizinhos, compara os valores de $k$ e imprime a matriz de confusão.

\begin{lstlisting}
python knn.py
\end{lstlisting}

{\sloppy Para demonstrar outra instância ou outro número de vizinhos, use os argumentos \mbox{\textbf{-{}-exemplo}} e \mbox{\textbf{-{}-k}}.\par}

\end{document}
